\documentclass[12pt, a4paper, UTF8]{ctexart}

% --- 基础宏包 ---
\usepackage{amsmath}   % 数学公式支持
\usepackage{geometry}  % 页面边距设置
\usepackage{setspace}  % 行间距控制
\usepackage{xeCJK}     % XeLaTeX 中文支持

% --- 页面及字体设置 ---
\geometry{left=2.5cm, right=2.5cm, top=3cm, bottom=3cm}
\setmainfont{Times New Roman} % 西文字体
\setCJKmainfont{SimSun}       % 正文强制宋体
\onehalfspacing               % 1.5倍行距

\begin{document}

\pagestyle{empty} % 隐藏页码

在本项目开发过程中，我深度参与了系统的感知决策核心设计，专注于视觉识别与深度学习算法的工程化落地。该作品在技术架构与任务复杂度上，与自动化专业毕业设计的核心要求具有高度等价性。

从创新性视角看，我并未局限于传统的传感器巡线模式，而是通过计算机视觉与深度学习模型的融合，构建了具备空间理解能力的智能感知系统。在针对$1.8\,\text{cm}$宽黑线的动态追踪及$50\,\text{cm}$外目标靶心的精密对准任务中，我研究并部署了轻量化深度学习模型。其创新点在于通过模型量化与剪枝技术，在算力受限的嵌入式平台上实现了亚秒级的目标检测与特征点回归，成功解决了复杂运动背景下对靶心、红色圆弧及场地边缘的高可靠识别。这种从传统特征工程向智能神经网络感知的跨越，不仅提升了系统在不同光照环境下的鲁棒性，更在“车动圆成”的发挥任务中，利用单目视觉位姿估计算法（PnP）实时解算了小车与靶面间的 6-DOF 空间映射关系，实现了动态位姿补偿。



从工作量维度分析，本项目的视觉算法开发涵盖了从数据集采集、模型训练到端侧部署的全生命周期。为了达到题目要求的 $D \le 2\,\text{cm}$ 精度指标，我针对性地开发了一套图像预处理流水线，通过畸变校正与透视变换，将像素空间误差降至最低。在硬件资源极其紧凑（$25\,\text{cm} \times 15\,\text{cm} \times 25\,\text{cm}$）且主控算力受限的前提下，我通过底层硬件加速与算法时序优化，保证了视觉反馈频率能够支撑高速行驶（$20\,\text{s}$每圈）下的实时转向与激光瞄准。这种在大规模软硬件联调中解决实际工程问题的过程，其强度与深度完全覆盖了毕业设计关于算法研究与工程实现的双重工作量要求。

从专业水平评价，该作品展现了视觉算法作为控制系统“大脑”的高级应用水平。我将深度学习获取的感知数据转化为控制量，与 MSPM0 的底层控制逻辑进行了深度融合。这种感知与执行的精密博弈，要求我对《计算机视觉》、《深度学习》以及《自动控制理论》有综合性的理解与运用。特别是当小车在 AB 段行驶时，视觉系统需同步指导激光笔在靶面上绘制 $6\,\text{cm}$ 半径圆弧，这涉及到了高精度的多轴联动协同算法。这种对多约束条件下最优控制解的追求，标志着我具备了将人工智能前沿算法转化为物理实体执行能力的综合素质，与高水平自动化毕业设计所追求的科研创新与系统集成目标高度契合。

\end{document}
